\begin{itemize}
    \item[а)] Пусть $g \in Aut(K)$. Какие значения может принимать $g(\sqrt[3]{2})$?
    \par \Proof Так как $(\sqrt[3]{2})^3-2=0 \Rightarrow g^3(\sqrt[3]{2})=2 \Rightarrow g$ переводит $\sqrt[3]{2}$ в корни многочлена $x^3-2$. Тогда $g(\sqrt[3]{2}) \in \{\sqrt[3]{2}, \omega \sqrt[3]{2}, \omega^2 \sqrt[3]{2}\}$. \EndProof
    
    \item[б)] Элемент $Aut(K)$ однозначно задаётся перестановкой корней многочлена $f(x)$
    \par \Proof Корня три $x_1, x_2, x_3$, а значит перестановок 6, а значит если задать перестановку корней $\sigma$, то получим, что $g(x_i)=g(\sigma(x_i))$, а это автоморфизм (проверить вручную нетрудно). Тогда по перестановке можно однозначно задать элемент из $Aut(K)$. \EndProof
    
    \item[в)] Группа $Aut(K)$ изоморфна $S_3$
    \par \Proof Так как данное нам поле — поле разложения многочлена $x^3-2$, расширение нормально. Тогда с учетом того, что степень расширения 6, количество автоморфизмов тоже 6. Тогда группа Галуа и есть $S_3$. \EndProof

    \item[г)] Рассмотрим подгруппу $A_3 \subset S_3$. Найдите поле $K^{A_3}$
    \par \Proof Рассмотрим $g \in Aut(K)$, которая соответствует циклу $(123)$. Тогда $g(\omega)=g(\frac{x_2}{x_1})=\frac{g(x_2)}{g(x_1)}=\frac{x_3}{x_2}=\omega$. Из соображения размерности получим, что $K^{A_3}=\Q(\omega)$. \EndProof
    
    \item[д)] Рассмотрим три подгруппы $H$, каждая из которых порождена одной транспозицией. Найдите $K^H$ для каждой такой подгруппы
    \par \Proof Проделав то же самое для перестановок (132, 213, 321) можно получить, что $\omega$ ни во что хорошее не перейдет (как и любая линейная комбинация), а значит $K^H=\Q$(корень который не участвует в транспозиции). \EndProof
    
    \item[е)] Расширение $K$ имеет четыре нетривиальных подполя, причём эти поля соответствуют подгруппам $S_3$. Более того, нормальным подгруппам $H \subset S_3$ соответствуют нормальные расширения $K^H \supset \Q$.
    \par \Proof Нетривиальные подгруппы из $S_3$ : $(123),(23),(12),(13)$. Тогда для $(23)$ получим, что $\omega \rightarrow \overline{\omega}$. И для нее $K^H=\Q(\sqrt[3]{2})$. Аналогично, рассмотрев подгруппу из $(x_ix_j), K^H=\Q(x_k)$, где $i,j,k$ — различные элементы из $1, 2, 3$. Тогда мы описали для нормальных подгрупп из двух элементов все, а для (123) описано в пункте г). \EndProof 
\end{itemize}